\documentclass[10pt,a4paper]{amsart}
\usepackage[margin=19mm]{geometry}
\usepackage{amsmath,amssymb,mathtools}
\usepackage[scr=rsfs,cal=euler]{mathalfa}
\usepackage{xfrac}
\usepackage[unicode,hidelinks,bookmarks=false]{hyperref}
\newcommand{\ie}{i.e.\ }
\newcommand{\cf}{cf.\ }
\newcommand{\resp}{respectively\ }
\newcommand{\Subsection}{\subsection}
\providecommand{\nobreakdash}{\nobreak\hbox{-}}
\setlength{\emergencystretch}{2em}
\multlinegap0pt
\newtheorem{proposition}{Proposition}
\renewcommand{\AA}{\mathbf{A}}
\newcommand{\FF}{\mathbf{F}}
\newcommand{\ZZ}{\mathbf{Z}}
\title{Appendix to ``Maximal curves of genus $5$ over finite fields''}
\author{Leolin Nkuete, Antigona Pajaziti, Hamide Suluyer, Rabia G\"ul\c{s}ah Uysal}
\date{Source: arXiv:2509.14871v2 (CC BY 4.0)}
\begin{document}
\maketitle

\noindent\emph{Source and licence.} Appendix to ``Maximal curves of genus $5$ over finite fields''. Version arXiv:2509.14871v2 (CC BY 4.0), distributed by arXiv under the Creative Commons Attribution 4.0 International licence, http://creativecommons.org/licenses/by/4.0/, as recorded in the arXiv licence metadata for this exact version and shown on the abstract page https://arxiv.org/abs/2509.14871v2. Full licence notice: \texttt{licenses/arxiv-2509.14871-CC-BY-4.0.txt}.\par\smallskip

\noindent\textit{Editorial scope.} Counted unit: the article's proposition proposition (source0.tex lines 114--116). The article's complete original proof of it is delivered with it (source0.tex lines 118--168, one \texttt{proof} environment of 51 lines). No mathematics is written, repaired or re-derived: the statement and the proof are the article's own lines, in the article's own notation, and no retained line is substituted or re-flowed.  No further source-local context is needed: every cross-reference of every retained line resolves inside the two delivered spans.  Nothing was omitted from any retained span, so the package carries no omission record.  No retained line contains a citation, so the wrapper carries no bibliography and no bibliography entry.  No retained line contains a figure environment, an image-inclusion command or drawing code, so no image is delivered and the package declares no asset dependency.  Editorial formatting only, in addition: an AMS \texttt{amsart} preamble that restates the article's own declarations and macros used by the retained lines, namely the environments \texttt{proposition} on separate counters, together with the standard packages those lines need.  The retained environments print in this document as Proposition 1 in order of appearance, whereas the article prints the same items with the numbering of its own full document.  The complete native source of the article as distributed in the arXiv e-print for this exact version is bundled unchanged as \texttt{sources/185276/source0.tex}.
\begin{proposition}
If $q=5$, then the curve $E$ cannot be isomorphic to $E'$.
\end{proposition}

\begin{proof}
If $q=5$, the degree $2$ divisor $D$ can only be $2P$ for some
$\FF_q$-rational point $P$ of $E$. In this case the order of
$U/\FF_q^*U_D$ is equal to $q=5$. The exact sequence $(\star)$ modulo fifth powers
leads to the exact sequence
\[
U/\FF_5^*U_D
\longrightarrow
\AA_E^*/{\AA_E^*}^5K^*U_D
\longrightarrow
\AA_E^*/{\AA_E^*}^5K^*U
\longrightarrow 0.
\]
If $E$ were isomorphic to the elliptic curve $E'$ given by
$y^2=x^3-2x+2$, the leftmost homomorphism is zero. In other words, $U$ is
contained in the subgroup ${\AA_E^*}^5K^*U_D$ of $\AA_E^*$. To see this, we first observe that the group $E(\FF_5)$ acts transitively on
itself by translations. Therefore we may take for $P$ any of the five points in
$E(\FF_5)$. We choose $P=(1,1)$. Then we put $Q=(2,-1)$ and let
$\xi \in \AA_E^*$ denote an id\`ele whose image in $\mathrm{Div}(E)$ is the
divisor $Q-O$, where $O$ is the point of $E$ at infinity. It is not difficult to see that the divisor of the function
\[
h =
\frac{(y+1)^2(x-2)}{y-2x-2}
\]
is equal to $5(Q-O)$. Therefore both $\xi^5$ and $h$ are elements of
$\AA_E^*$ whose images in $\mathrm{Div}(E)$ are equal to $5(Q-O)$. It follows
that $v = \frac{\xi^5}{h}$ is in $U$. Since $h(P)=2$ and
\[
h-2 =
\frac{(x-1)(2y+x^3-x^2+2x)}{y-2x-2}
\]
has a simple zero in $P$, the id\`ele unit $v$ generates $U/U_D$. Therefore
every $u \in U$ can be written as $v^k w$ for some $k \in \ZZ$ and
$w \in U_D$. Since $v$ is in ${\AA_E^*}^5K^*$, this proves that we indeed have $U \subset {\AA_E^*}^5K^*U_D$.
 

It follows that the natural map
\[
\AA_E^*/{\AA_E^*}^5K^*U_D
\mathop{\longrightarrow}^{\cong}
\AA_E^*/{\AA_E^*}^5K^*U
\]
is an isomorphism. By class field theory, the group on the left is isomorphic
to the Galois group of the maximal abelian exponent $5$ extension $L$ of $K$ of
conductor at most $D=2P$. Note that the function field $\FF_5(X)$ of the curve
$X$ is a subfield of $L$. Since $\AA_E^*/{\AA_E^*}^5K^*U_D$ is isomorphic to  $\AA_E^*/{\AA_E^*}^5K^*U$,
 the extension $K \subset L$ is unramified at \textit{all} places. Since we have $K \subset \FF_5(X) \subset L$,
 the same is true for $\FF_5(X)$. However, the cover $X \rightarrow E'$ is
actually ramified and hence we obtain a contradiction.
This proves the proposition.
\end{proof}

\end{document}
